\section*{Hoofdstuk \ref{ch:g12:proton-nmr} --- Proton-NMR}

\begin{solution}{exo:g12:proton-nmr:1}
Methaan: 1. Ethaan: 1 (de twee \ce{CH3} zijn gelijk). Propaan: 2 (twee
gelijke \ce{CH3}, één \ce{CH2}). Methanol: 2 (\ce{CH3} en \ce{OH}).
Methoxymethaan: 1.
\end{solution}

\begin{solution}{exo:g12:proton-nmr:2}
2 buren: een triplet, 1:2:1. 3 buren: een kwartet, 1:3:3:1.
0 buren: een singlet.
\end{solution}

\begin{solution}{exo:g12:proton-nmr:3}
In totaal \qty{40}{mm} voor 8 \omterm{def:g9:inside-the-atom:proton}{protonen}: \qty{5}{mm} per \omterm{def:g9:inside-the-atom:proton}{proton}. De
signalen stellen 3, 2 en 3 \omterm{def:g9:inside-the-atom:proton}{protonen} voor.
\end{solution}

\begin{solution}{exo:g12:proton-nmr:4}
\omterm{def:g11:functional-groups:carbonyl}{Aldehyde}: 9.7 tot \qty{10.0}{ppm}. Een \ce{CH3} in een keten: 0.7 tot
\qty{1.3}{ppm}.
\end{solution}

\begin{solution}{exo:g12:proton-nmr:5}
Zijn twaalf \omterm{def:g12:proton-nmr:equivalent}{equivalente protonen} geven één scherpe, intense lijn, die
gemakkelijk op nul te zetten is. Het gedeutereerde \omterm{def:g6:solutions-and-solubility:solute}{oplosmiddel} geeft zelf
geen protonsignaal, dat anders de signalen van de paar milligram monster
zou overspoelen.
\end{solution}

\begin{solution}{exo:g12:proton-nmr:6}
Twee signalen. \ce{CH2} naast chloor: 2 \omterm{def:g9:inside-the-atom:proton}{protonen}, een kwartet (3
buren), in het gebied 2.5--\qty{4.0}{ppm}. \ce{CH3}: 3 \omterm{def:g9:inside-the-atom:proton}{protonen}, een
triplet (2 buren), bij het gebruikelijke \ce{CH3}-gebied, iets hoger
omdat het chlooratoom twee \omterm{def:g7:atoms-and-molecules:atom}{atomen} verderop zit.
\end{solution}

\begin{solution}{exo:g12:proton-nmr:7}
Drie signalen: de twee \ce{CH3}, 6 \omterm{def:g12:proton-nmr:equivalent}{equivalente protonen}, een doublet (1
buur, de \ce{CH}); de \ce{CH}, 1 \omterm{def:g9:inside-the-atom:proton}{proton}, door 6 buren gesplitst
in 7 lijnen, in het \ce{H-C-O}-gebied (3.3--\qty{4.5}{ppm}); de
\ce{OH}, 1 \omterm{def:g9:inside-the-atom:proton}{proton}, een singlet.
\end{solution}

\begin{solution}{exo:g12:proton-nmr:8}
(a) propanon; (b) ethylethanoaat; (c) propaanzuur.
\end{solution}

\begin{solution}{exo:g12:proton-nmr:9}
De \ce{CH2} is gebonden aan het zuurstofatoom van de \omterm{def:g11:functional-groups:carboxylic-acid}{ester}, dat
elektronegatief is en zijn \omterm{def:g9:inside-the-atom:proton}{protonen} ontschermt: \ce{H-C-O}-gebied,
3.3--\qty{4.5}{ppm}. De \ce{CH3} zit één koolstofatoom verder en blijft
in het ketengebied.
\end{solution}

\begin{solution}{exo:g12:proton-nmr:10}
Propanon: een \ce{C=O} (\omterm{def:g11:functional-groups:carbonyl}{keton}, \qty{1715}{cm^{-1}}) en zes \omterm{def:g12:proton-nmr:equivalent}{equivalente
protonen}, één singlet. Propanal zou drie signalen geven, waaronder dat
van het aldehydeproton van \ce{CH3-CH2-CHO}, bij 9.7--\qty{10.0}{ppm}.
\end{solution}

\begin{solution}{exo:g12:proton-nmr:11}
Het \ce{OH}-proton wordt snel tussen \omterm{def:g7:atoms-and-molecules:molecule}{moleculen} uitgewisseld: het geeft
een singlet en splitst zijn buren niet. De \ce{CH2} wordt dus alleen
gesplitst door de 3 \omterm{def:g9:inside-the-atom:proton}{protonen} van de \ce{CH3}: een kwartet.
\end{solution}

\begin{solution}{exo:g12:proton-nmr:12}
Butaanzuur \ce{CH3-CH2-CH2-COOH}: 4 signalen, een triplet (\ce{CH3}), een
signaal van 6 lijnen (middelste \ce{CH2}, 5 buren), een triplet
(\ce{CH2-C=O}), een singlet bij 11--\qty{12}{ppm}. Ethylethanoaat: een
singlet, een kwartet, een triplet. Methylpropanoaat \ce{CH3-CH2-CO-O-CH3}:
een triplet, een kwartet (\ce{CH2-C=O}, 2.0--\qty{2.4}{ppm}), een singlet
(\ce{O-CH3}). Propylmethanoaat \ce{H-CO-O-CH2-CH2-CH3}: 4 signalen, een
singlet ver naar links (de H aan het koolstofatoom van de \ce{C=O}), een
triplet (\ce{O-CH2}), een signaal van 6 lijnen, een triplet (\ce{CH3}).
\end{solution}

\begin{solution}{exo:g12:proton-nmr:13}
Het \ce{OH}-signaal: de \ce{O-H}-protonen wisselen uit met het deuterium
van het zware water, en \ce{O-D} geeft geen protonsignaal. Een signaal
dat verdwijnt na schudden met \ce{D2O}, hoort bij een \omterm{def:g9:inside-the-atom:proton}{proton} aan O of N.
\end{solution}

\begin{solution}{exo:g12:proton-nmr:14}
Ethanol voegt een kwartet toe bij ongeveer \qty{3.69}{ppm}, een triplet
bij ongeveer \qty{1.23}{ppm} en een \ce{OH}-singlet. Het singlet van de
\omterm{def:g11:functional-groups:carboxylic-acid}{ester}: \qty{30}{mm} voor 3 \omterm{def:g9:inside-the-atom:proton}{protonen}, \qty{10}{mm} per \omterm{def:g9:inside-the-atom:proton}{proton} \omterm{def:g11:functional-groups:carboxylic-acid}{ester}. Het
triplet van ethanol: \qty{3}{mm} voor 3 \omterm{def:g9:inside-the-atom:proton}{protonen}, \qty{1}{mm} per \omterm{def:g9:inside-the-atom:proton}{proton}
ethanol. Omdat elk \omterm{def:g7:atoms-and-molecules:molecule}{molecuul} één \ce{CH3} aan deze signalen bijdraagt,
staan de hoeveelheden in de verhouding $1/10$: ongeveer één \omterm{def:g7:atoms-and-molecules:molecule}{molecuul}
ethanol op tien \omterm{def:g7:atoms-and-molecules:molecule}{moleculen} \omterm{def:g11:functional-groups:carboxylic-acid}{ester}.
\end{solution}

\begin{solution}{exo:g12:proton-nmr:15}
Vier signalen: de twee \ce{CH3} (6 \omterm{def:g9:inside-the-atom:proton}{protonen}) een doublet; de \ce{CH} (1
\omterm{def:g9:inside-the-atom:proton}{proton}), door 6 + 2 = 8 buren gesplitst in 9 lijnen; de \ce{CH2} (2
\omterm{def:g9:inside-the-atom:proton}{protonen}) naast de \ce{O}, een doublet, in het \ce{H-C-O}-gebied; de
\ce{OH} (1 \omterm{def:g9:inside-the-atom:proton}{proton}) een singlet. Het \ce{CH}-signaal heeft de meeste
lijnen.
\end{solution}

\begin{solution}{pb:g12:proton-nmr:1}
\textbf{1.} Een \omterm{def:g11:organic-skeletons:alkane}{alkaan} met vier koolstofatomen is \ce{C4H10}; één
\ce{C=O} haalt er twee waterstofatomen af en de twee zuurstofatomen
voegen er geen toe: \ce{C4H8O2}, de formule van zuren en \omterm{def:g11:functional-groups:carboxylic-acid}{esters} met vier
koolstofatomen.

\textbf{2.} Butaanzuur \ce{CH3-CH2-CH2-COOH}; ethylethanoaat
\ce{CH3-CO-O-CH2-CH3}; methylpropanoaat \ce{CH3-CH2-CO-O-CH3};
propylmethanoaat \ce{H-CO-O-CH2-CH2-CH3}.

\textbf{3.} De drie \omterm{def:g11:functional-groups:carboxylic-acid}{esters} hebben de estergroep gemeen; alle vier hebben
een \ce{C=O}.

\textbf{4.} Een \ce{C=O}, in het estergebied (bij ongeveer
\qty{1735}{cm^{-1}}).

\textbf{5.} De zeer brede \ce{O-H}-band van 2500 tot
\qty{3300}{cm^{-1}}: die ontbreekt, dus geen butaanzuur.

\textbf{6.} Nee: de drie \omterm{def:g11:functional-groups:carboxylic-acid}{esters} hebben dezelfde groepen.

\textbf{7.} Drie.

\textbf{8.} Een ethylgroep \ce{CH3-CH2-}: de \ce{CH2} gesplitst door 3
\omterm{def:g9:inside-the-atom:proton}{protonen} (kwartet), de \ce{CH3} door 2 (triplet).

\textbf{9.} Er zitten geen \omterm{def:g9:inside-the-atom:proton}{protonen} aan de naburige \omterm{def:g7:atoms-and-molecules:atom}{atomen}: een
\ce{CH3} gebonden aan het koolstofatoom van de \ce{C=O} of aan het
zuurstofatoom van de estergroep.

\textbf{10.} De \ce{CH2} is aan zuurstof gebonden (\ce{H-C-O}-gebied).

\textbf{11.} Een triplet (\ce{CH3}), een kwartet bij 2.0--\qty{2.4}{ppm}
(\ce{CH2} naast \ce{C=O}) en een singlet in het \ce{H-C-O}-gebied
(\ce{O-CH3}). Het kwartet zou ver onder \qty{4.12}{ppm} liggen en het
singlet ver boven \qty{2.04}{ppm}: niet de onbekende stof.

\textbf{12.} Vier signalen (een singlet ver naar links, een triplet, een
signaal van 6 lijnen, een triplet): niet de onbekende stof, die er drie
heeft.

\textbf{13.} Ethylethanoaat, \ce{CH3-CO-O-CH2-CH3}: singlet
\qty{2.04}{ppm} = \ce{CH3-C=O}; kwartet \qty{4.12}{ppm} = \ce{O-CH2};
triplet \qty{1.26}{ppm} = de \ce{CH3} aan het eind.

\textbf{14.} Singlet 3, kwartet 2, triplet 3.

\textbf{15.} $72 / 8 = \qty{9}{mm}$ per \omterm{def:g9:inside-the-atom:proton}{proton}.

\textbf{16.} Singlet $3 \times 9 = \qty{27}{mm}$; triplet
\qty{27}{mm}.

\textbf{17.} $18 + 27 + 27 = \qty{72}{mm}$.

\textbf{18.} Ja: kleine \omterm{def:g11:functional-groups:carboxylic-acid}{esters} staan bekend om hun fruitige geur.

\textbf{19.} De trede van het kwartet is \qty{18}{mm} van de
\qty{72}{mm}.
\end{solution}
